\chapter{Synthetic Nervous System Simulation}\label{ch:sns}
\graphicspath{{figs/60-sns/}{figs/60-sns/AnimatLab_reference/}}

The actuator and torque models of Chapters~\ref{ch:muscles} and~\ref{ch:torque} describe the robot's mechanics. Studying their interaction with neural control also requires a simulation that connects the mechanical model to a synthetic nervous system (SNS). I first attempted this integration in AnimatLab by combining a published walker with a spinal controller. The model produced stepping-like motion in air, but hand tuning did not produce walking. This chapter explains how that attempt motivated the development of software tools for the assembly, testing, and comparison of neuromechanical models.

\section{The AnimatLab Bipedal Walker}\label{sec:prior_walker}

For a course project in ME 610: Neuromechanical Modeling, Morrow and I combined the bipedal mechanics of \citet{li_bipedal_2017} with the two-level central pattern generator (CPG) of \citet{deng_neuromechanical_2019} in AnimatLab \citep{cofer_animatlab_2010}. The mechanical model had a pelvis and bilateral thigh, shank, foot, and toe segments, with sagittal-plane joints and antagonist Hill-type muscles at the hip, knee, and ankle. I added a passive toe joint. The initial controller used one rhythm generator (RG) to control both sides, with a separate pattern-formation layer. We added afferent feedback at the MN level. Figure~\ref{fig:animatlab_model_architecture} shows the model and controller organization.

We adjusted applied currents, synaptic weights, and muscle parameters by hand. The suspended model developed alternating muscle activity and stepping-like motion, presented at the 2021 NIH BRAIN Initiative Investigators Meeting \citep{bolenTorqueValuesArtificialKnee2021June15-17}. It did not achieve weight-bearing walking.

I subsequently expanded the neural model into left- and right-side circuits with interconnected rhythm generators (RGs). I constructed the right-side subsystems from the left-side organization and linked the two RGs through commissural pathways. I added afferent feedback to the RG and PF layers, and heel- and toe-contact signals during the subsequent development. This later model is referred to as the Walker Two-Layer (W2L) model. Further circuit schematics and implementation records are provided in Appendix~\ref{app:simulation_verification}.

The AnimatLab model established that the combined circuit could generate alternating activity and move the suspended limbs. To investigate the contribution of sensory feedback and actuator mechanics, I needed to vary those components while retaining a specified neural organization and mechanical support condition. Extending that approach to the human-scale Gait2392 model and its BPA-routed robot counterpart also required checking joint directions, muscle paths, actuator forces, and the signals exchanged with the controller.

\begin{figure}[p]
\centering
\makebox[\textwidth][l]{\textbf{(A)}}\par\smallskip
\includegraphics[width=0.64\textwidth,height=0.21\textheight,keepaspectratio]{OurWalker.png}
\par\medskip
\makebox[\textwidth][l]{\textbf{(B)}}\par\smallskip
\includegraphics[width=\textwidth,height=0.50\textheight,keepaspectratio]{connor_ben_shared_rg_original.png}
\caption[Original AnimatLab walker and shared-RG controller]{Original AnimatLab model developed by Morrow and Bolen: (A) mechanical walker and (B) neural controller. The shared four-neuron RG drives bilateral PF groups, with a hip group and a knee/ankle group for each leg. This is the original controller, before the later split into one RG per side and the addition of feedback to the RG and PF layers and foot-contact feedback.}
\label{fig:animatlab_model_architecture}
\end{figure}

\section{Knowledge Gaps and the New Tools}\label{sec:sns_gaps}

Researchers studying neuromechanical control need to distinguish movement generated by a neural circuit from movement shaped by sensory feedback, actuator mechanics, and external support. For the humanoid models studied here, the unresolved questions were whether a specified circuit retained its behavior when implemented in a different simulator, and how its output changed when coupled to human musculotendon mechanics or artificial muscles. A successful neural trace alone could not answer the mechanical question, and a moving model alone could not establish that the intended neural pathways caused the movement.

I addressed this gap by constructing models in which the neural pathways, actuator law, sensory inputs, and support conditions can be specified and varied, together with tests that distinguish circuit behavior from implementation and coupling errors. The resulting contributions are:
\begin{enumerate}
    \item A custom Li--Deng walker and controller, developed in AnimatLab and implemented in MuJoCo with SNS-Toolbox, with matched implementations of its contact-driven neural core in NumPy and Simulink.
    \item Gait2392-based neuromechanical models with alternative neural controllers, and a robot counterpart with BPA routing and the actuator-force model from Chapter~\ref{ch:muscles}.
    \item The SNS-Simscape block library and GUI for block programming neural pathways, including sensory, neuron, synapse, activation, and custom BPA-force blocks.
    \item A SolidWorks-derived robot assembly in Simscape Multibody coupled to a neural controller and BPA actuators using my force model.
    \item Integration of SNS-Simscape controllers with MuJoCo through the MathWorks bidirectional bridge, together with model-import procedures for musculoskeletal and CAD-derived mechanisms.
    \item A five-stage curriculum and experiment workflow for parameter searches, adjustable support, neural stimuli, pathway interventions, mechanical perturbations, and recording neural and mechanical outputs.
    \item Implementations of literature-derived circuits and selected spiking-neuron counterparts for comparing neural mechanisms and coding assumptions.
\end{enumerate}
Table~\ref{tab:sns_capability} gives selected functional-verification tests and proof-of-concept demonstrations. These establish that the tools can construct, exercise, and compare the specified models. The feedback experiments also provide a limited mechanistic result: removing particular sensory pathways can eliminate stepping in a fixed supported configuration, identifying a dependence that would not be apparent from its intact motion alone.

Existing software supplied the underlying simulation capabilities. OpenSim provided musculoskeletal geometry, MyoConverter translated models to MuJoCo, and SNS-Toolbox provided neuron and synapse formulations \citep{delp_opensim_2007, caggiano_myosuite_2022, todorov_mujoco_2012, nourse_sns_2023}. Simulink provided graphical block programming, and Simscape Multibody provided mechanical modeling and import of SolidWorks-derived mechanisms. My contributions are the particular model implementations, custom blocks, actuator integrations, and evaluation procedures built with those packages. The SNS-Simscape library makes the neural pathways explicit in a block diagram, while the MuJoCo implementations support repeated scripted evaluations and parameter searches. Matched-input tests check selected circuits before their outputs are interpreted as neural control of a mechanical model.

\section{Software Tools}\label{sec:preliminary_sim_methods}

The tools provide three complementary capabilities: constructing a neuromechanical model, testing its component behavior and signal exchange, and evaluating controlled changes in neural or mechanical conditions. Table~\ref{tab:sns_capability} summarizes representative demonstrations. Figure~\ref{fig:pipeline} places the simulation tools within the longer-term robot-control workflow.

\begin{table}[htbp]
\caption[Tools developed and representative demonstrations]{Tools developed in this project and representative functional-verification tests and proof-of-concept demonstrations. The examples illustrate the toolset rather than exhaustively evaluating every implementation; each result refers to the named circuit or mechanical model.}
\label{tab:sns_capability}
\small
\begin{tabularx}{\textwidth}{@{}>{\raggedright\arraybackslash}p{0.23\textwidth}*{2}{>{\raggedright\arraybackslash}X}@{}}
\toprule
\textbf{Tool developed} & \textbf{Function} & \textbf{Representative demonstration} \\
\midrule
Li--Deng/W2L model implementations & Combine the walker mechanics and neural architecture; transfer selected controller components between environments & Matched contact-driven rhythm cores reproduce \qty{1.000}{\second} periods with trace correlation at least 0.99996 \\
Gait2392/SNS model implementations & Evaluate alternative controllers on the human reference and custom actuator/routing models & Individual-actuator tests check torque directions; selected human-model controllers produce bilateral stepping-like activity under support \\
SNS-Simscape library and GUI & Construct neural pathways graphically and connect sensory, neural, and actuator blocks & The loaded-arm reflex demo gives \ang{2.00} sag with feedback and \ang{7.85} without it; actuator blocks agree with reference force calculations \\
W2L BPA-force integration & Apply the Festo force model within the coupled simulation & The 12-actuator BPA model demonstrates air stepping and retained-support standing, with peak pressure \qty{372}{\kPa} \\
SNS-Simscape/MuJoCo integration & Exchange activation commands and mechanical measurements through the MathWorks bridge & Single-muscle and control-rate tests reproduce the Python-reference sensor readback over the scored windows \\
Curriculum and experiment drivers & Automate staged tuning, support changes, neural stimuli, and mechanical perturbations & Saved stage-specific trials, candidate replays, and fixed-controller pathway and body-change comparisons demonstrate the workflow \\
Literature-circuit and spiking implementations & Test selected published neural behaviors and compare graded and spiking formulations & The reduced Rybak circuit passes 13 drive/deletion checks; selected spiking counterparts have paired implementation tests \\
\bottomrule
\end{tabularx}
\end{table}

\subsection{Model Conversion and Actuator Coupling}\label{sec:sns_existing_tools}

My first conversion attempt used the robot model, which combined BPA routing on the Gait2392 skeleton with unchanged human muscle parameters. MyoConverter's muscle-equilibrium checks rejected infeasible combinations of routed length and muscle parameters. I therefore separated conversion of the human reference from reconstruction of the robot routing. This distinction matters because an error caused by an artificial-actuator route should not be interpreted as a property of the original human model.

The human-model conversion also required specific repairs. Three small hip rotators per side, the quadratus femoris, gemelli, and piriformis, were disabled at model load. The converted model retains 92 actuator entries, but these six contribute no force. The model retains the OpenSim and International Society of Biomechanics coordinate convention; frame transformations and individual-actuator tests check that joint motion and torque directions have the intended signs.

The source represents patella motion through muscle insertion locations that vary with knee angle. In the converted model, limits on the joints representing these moving path points conflicted with their coupling constraints. Removing those conflicting limits allowed the coupled path points to follow their prescribed trajectories. The rectus femoris path also required a patella-tracking via point to restore its knee-extension action. Explicit collision pairs introduced unwanted contacts, so contact was restricted to the foot--ground pairs. The anatomical joint limits also had to be re-enabled; these driver joints were treated separately from the coupled path-point joints. Passive joint torques supplemented the joint behavior. These repairs explain why conversion needed checks beyond successful loading: a model could load while producing an incorrect torque direction or unintended contact force.

Comparisons with the OpenSim attachment paths and single-actuator tests checked the transferred geometry and force directions. They do not establish that every muscle reproduces the source force throughout its range or that the controller produces walking. The neural and mechanical models exchange signals at a \qty{2}{\milli\second} solver time step. Appendix~\ref{sec:app_conversion_repairs} records the individual repairs and their scope.

For artificial-muscle simulations, I implemented an actuator-force calculation that evaluates the quasi-static BPA model from Chapter~\ref{ch:muscles} at the current simulated actuator length and commanded pressure. The calculated tensile force acts along the model's specified tendon route. Length, velocity, force, and foot-contact measurements are then returned to the neural controller through sensory channels. This closes the neural-command--actuator-force--sensory-feedback loop. The calculation represents the measured static force dependence; it does not add valve-flow or pressure-transient dynamics. Routing geometry and dynamic force transmission therefore require their own checks beyond the isometric force fit.

\subsection{Block Programming Neural Pathways}\label{sec:sns_new_tools}

The SNS-Simscape library implements neuron, synapse, afferent-encoder, activation, and muscle-force elements as editable Simulink blocks. The library retains the conductance-based neuron and synapse formulation supplied by SNS-Toolbox; it does not introduce a new neuron model. Verification uses selected circuits under matched inputs. Figure~\ref{fig:simulink_sns_implementation} shows the core blocks and graded-neuron signal flow.

I imported OpenSim-derived and SolidWorks-derived mechanisms into Simscape Multibody and coupled the SolidWorks robot assembly to a neural controller and BPA actuators using my force model. The block library also supports reduced-order demonstrations in which the sensory pathway and mechanical response can be inspected directly. These smaller models provide interpretable checks of the components used in larger assemblies.

\begin{figure}[htbp]
\centering
\includegraphics[width=\textwidth]{pipeline.pdf}
\caption[Simulation-to-robot pipeline]{Model conversion, closed-loop simulation, and proposed hardware deployment. Solid simulation interfaces connect neural activation to BPA force and return afferent signals. The embedded valve-control and network-evaluation system is future work (Chapter~\ref{ch:futurework}).}
\label{fig:pipeline}
\end{figure}

\begin{figure}[p]
\centering
\includegraphics[width=0.94\textwidth,height=0.42\textheight,keepaspectratio]{sns_library_overview.pdf}
\par\medskip
\includegraphics[width=0.94\textwidth]{simulink_neuron_detail.pdf}
\caption[Simulink SNS library blocks and neuron implementation]{Core SNS-Simscape blocks (top) and graded-neuron implementation (bottom). Summed currents drive membrane integration; threshold, inverse-slope gain, and saturation bound the output. Voltage feeds both current calculations. The generic BPA demonstration remains provisional.}
\label{fig:simulink_sns_implementation}
\end{figure}

\section{Verification of the Implemented Interfaces}\label{sec:preliminary_sim_results}

The tests address three distinct questions: whether a circuit produces its expected response, whether implementations agree under matched inputs, and whether the neural and mechanical models exchange signals in a closed loop.  The larger evaluation sets and their limitations are collected in Appendix~\ref{app:simulation_verification}; neuron equations and calibration details remain in Appendix~\ref{app:neuron_equations}.

\subsection{Reflex Response}\label{sec:sns_simulink_demos}

The library assembles Ia stretch excitation, reciprocal inhibition, and Ib autogenic inhibition in a reduced-order antagonist knee model (Figure~\ref{fig:simulink_reflex_topology}). A second demonstration tests feedback in an arm holding a beverage container as it is loaded. With the reflex enabled, arm sag was \ang{2.00}, compared with \ang{7.85} when it was disabled. The reduced sag demonstrates that the implemented sensory pathway changes motor activation in a direction that opposes the applied load. It is a functional test of the coupled neural and actuator blocks, rather than an experimental measurement of human-arm behavior.

\begin{figure}[p]
\centering
\includegraphics[width=\textwidth]{KneeReflex_circuit.pdf}
\caption[Reduced-order knee-reflex circuit topology]{Reduced-order antagonist knee-reflex circuit. Triangles mark excitation, and filled circles mark inhibition. Ia input excites the homonymous motor pool and inhibits its antagonist; Ib input inhibits the homonymous pool. Motor output drives activation and BPA-force blocks. Mechanical parameters and the generic BPA block remain provisional.}
\label{fig:simulink_reflex_topology}
\end{figure}

\subsection{Matched Neural Implementations}\label{sec:crossengine_results}

The contact-driven W2L rhythm core was implemented in NumPy and native Simulink blocks from a common connectivity specification. Both implementations received the same input protocol. Burst periods matched at \qty{1.000}{\second}, all ten expected bursts were reproduced, and voltage-trace correlation was at least 0.99996 (Figure~\ref{fig:w2l_crossengine}). The agreement is specific to this circuit and input protocol.

Removing the phasic heel-contact inputs caused the circuit to latch in both implementations. The alternating activity therefore depends on the specified sensory drive.

\begin{figure}[htbp]
\centering
\includegraphics[width=\textwidth,height=0.55\textheight,keepaspectratio]{sns_w2l_traces.png}
\caption[W2L rhythm-core traces across engines]{W2L rhythm-generator voltages under matched inputs in NumPy and Simulink. Burst periods and counts agree; full-trace correlation is at least 0.99996. The residual \qty{2}{\milli\second} lag is one discrete solver step. The test concerns the driven rhythm core, not walking mechanics.}
\label{fig:w2l_crossengine}
\end{figure}

\subsection{Closed-Loop Model Coupling}

The MathWorks MuJoCo--Simulink bridge exchanges activation commands and mechanical measurements between a Simulink controller and an existing MuJoCo mechanical model. I paired this bridge with SNS-Simscape and tested single-muscle activation, different control rates, and sensor readback against a Python integration. Readback was bit-exact over the scored window. An exported neural network also drove all 92 actuator channels during a \qty{5}{\second} run without afferent return; a separate closed-loop reflex experiment exercised both commands and sensory measurements. These tests verify signal exchange through the interface. Under tonic drive, one Simulink circuit realization reached a fixed point while its NumPy reference oscillated, showing why successful signal exchange must be followed by circuit-specific dynamical checks before transferring tuned parameters.

\subsection{Feedback Contribution at Fixed Controller Settings}\label{sec:sns_feedback_test}

At fixed controller and drive settings, removal of contact and Ib force-feedback pathways reduced stepping in the supported W2L model from 33 left and 28 right cycles to one left and no right cycles. Forward displacement decreased from \qty{1.16}{\meter} to \qty{0.09}{\meter}. These measures cover the \qty{25.499}{\second} scored window after settling within each \qty{30}{\second} simulation. The harness configuration was unchanged, although its measured support contribution differed between conditions. Ia and II pathways remained active. This comparison demonstrates the contribution of the removed pathways at this operating point; it uses the standard muscle-actuated W2L model.
\begin{figure}[t]
\centering
\includegraphics[width=0.85\textwidth]{ablation_pair_metrics_20261026.pdf}
\caption[Contact and force-feedback ablation]{Contact and Ib force-feedback ablation at fixed controller and drive settings in the supported W2L model. Counts and displacement refer to the post-settling scored window. The retained Ia and II pathways distinguish this test from complete deafferentation.}
\label{fig:sns_ablation_pair}
\end{figure}

\section{Discussion: What the Tools Establish}\label{sec:sns_discussion}

The principal contribution is a set of executable neuromechanical models and experimental tools that make the roles of neural pathways and actuator mechanics accessible to controlled tests. The initial AnimatLab model established the combined walker and CPG organization. The Gait2392 implementations extend controller development to a human-scale model with many joints and actuator channels, while the BPA implementations connect the measured actuator-force model to neural control. SNS-Simscape provides a graphical representation of the pathways, and the curriculum and experiment drivers provide repeatable changes in drive, support, feedback, and perturbation conditions.

The representative results answer different parts of the knowledge gap. Agreement between the NumPy and Simulink rhythm cores establishes that the specified circuit response can be preserved across those implementations under matched inputs. The loaded-arm comparison establishes a corrective effect of feedback in the coupled block model. The fixed-controller W2L ablation shows that its tested stepping behavior depends on contact and Ib-force feedback. The bridge tests establish the command and measurement exchange needed to carry out such experiments across software environments. Together, these results provide an operational method for separating implementation agreement, sensory contribution, and mechanical response.

The supported robot evaluations produced stepping-like motion and, in some configurations, forward progression, but no harness-free walking. Reducing tonic neural drive did not remove mechanical support. Spiking-neuron comparisons preserved a continuous motor-activation interface while changing rhythm timing and, in some runs, limb coordination (Appendix~\ref{app:simulation_verification}).

Researchers can use these tools to compare candidate neural pathways, actuator descriptions, and support conditions while recording their effects on the same defined outputs. The supported demonstrations establish this experimental capability; they do not establish balance without the support rig or transfer of every controller between models. Chapter~\ref{ch:futurework} discusses extensions to weight-bearing control and hardware. Those extensions build on the models and tests already demonstrated here.